% ============================================================================
% CH 2 — BACKGROUND AND RELATED WORK     [slugs: background + related-work]
% STATUS: MERGED from the two former chapters (Matteo's call, 2026-07-14).
% SQ-7: the DBIS sample skeleton separates Background and Related Work; the
% supervisor may ask for a split — note that it would be quite hard to
% untangle (positioning passes are interleaved with the teaching sections at
% paragraph level; the interleaving is the point of the merge).
% Discipline: teach-then-position along two axes; every positioning passage
% makes explicit relate-to-ours moves (DBIS norm: argumentative, NOT a
% survey); chapter ENDS with The Gap. Works we BUILD ON are taught at working
% depth; works we POSITION AGAINST get relate-to-ours passes only.
% The old rw-eval-explanations section relocated to Ch "Interpretability
% Quantified" (sec:iq-taxonomy) — it only serves that chapter.
% Material map: ../outline.md#background
% ============================================================================
\cleardoublepage
\chapter{Background and Related Work (12--15)}   % page goal — scaffolding, DELETE before submission
\label{ch:background}
\label{ch:related}   % kept alive so old cross-references keep resolving

% CHAPTER INTRO (bonbon, unnumbered) — absorbs the old "Scope: Two Axes"
% section. Three jobs:
%  1. Name the intersection: axis 1 = how attribute information enters an MF
%     backbone (the mechanism); axis 2 = how models earn their
%     interpretability (the goal). Each axis alone is well-populated; the
%     thesis lives at the crossing.
%  2. Draw the boundary — explicit exclusions with one-line reasons
%     (modality-light positioning, vault 2026-07-13_1817): review-corpus
%     explanation methods, image/multimodal and LLM-based explainers
%     (catalogue attributes only; H&M has no reviews); bias/fairness parked
%     by deliberate omission (Conclusion owns it).
%  3. Reading guide / division of labor: teaching sections vs positioning
%     passes, and the promise that the chapter closes with the gap.

\section{Matrix Factorization in Recommender Systems}
\label{sec:bg-mf}
\vault{2026-06-16_1631_protomf-item-ranking-basecase}
% Renamed from "MF for Implicit Feedback" (Matteo, 2026-07-14): opener must
% establish the task + model class; implicitness is a stage inside, not the
% section's identity. Build order:
%  1. The recommendation task: users x items, observed interactions.
%  2. MF: low-rank factorization, entities as latent factor vectors, score =
%     inner product (the invariant the whole lineage keeps). Frame MF as the
%     workhorse of collaborative filtering (CF as the family; sets up the
%     "collaborative = what trains, content = how parameterized" move later).
%  2b. RELEVANCE MOVE (Matteo, 2026-07-14): CF's importance TO THIS DAY —
%     not a historical technique: well-tuned MF/CF baselines remain
%     embarrassingly competitive with neural recommenders (Dacrema et al.
%     "Are we really making much progress?"; Rendle et al. 2020 "NCF vs MF
%     Revisited"), and CF backbones power production systems at scale.
%     Payoff sentence: making THIS model family interpretable is therefore
%     high-leverage — we harden the workhorse, not a boutique architecture.
%  3. The implicit-feedback setting (the ProtoMF base case): preference
%     INFERRED from behavior (buys, clicks), binary positive-only — vs
%     explicit ratings. Honest sentence: ML-1M is explicit data implicitized
%     (splitter keeps ratings >= 3.5, drops the value); H&M purchases are
%     genuinely implicit. One-class asymmetry: a missing interaction is
%     UNLABELED (dislike or never-seen), not an observed negative.
%  4. Two consequences: (a) negative sampling — training manufactures
%     presumed negatives from non-interactions; (b) the task is item RANKING,
%     not matrix completion — HR@K / NDCG@K (definitions live in the
%     eval-framework chapter).
% Still-MF argument + latent taxonomy: vault 2026-07-13_2033.
\begin{itemize}
    \item Still MF through the whole lineage: only the \emph{parameterization of the factors} changes. \vault{2026-07-13_2033_what-remains-latent-after-grounding-still-mf}
    \item Latent $\neq$ learned $\neq$ anchored: an anchored embedding has latent coordinates but an observed identity. \vault{2026-07-13_2033_what-remains-latent-after-grounding-still-mf}
    \item Classic MF answers ``what is latent?'' with \emph{everything}; our models make it an empirical, per-model number. \vault{2026-07-13_2033_what-remains-latent-after-grounding-still-mf}
    \item Latency is cornered and quantified, never eliminated — smallest at the noid--noid merge. \vault{2026-07-13_2033_what-remains-latent-after-grounding-still-mf}
\end{itemize}

\section{Interpretable Machine Learning}
\label{sec:bg-iml}
\label{sec:bg-iml-intrinsic}    % kept alive: old subsection labels parked on
\label{sec:bg-iml-concepts}     % the section until the structure is decided
\label{sec:bg-iml-additivity}
\label{sec:rw-prototypes}       % old RW sections fold in here (positioning)
\label{sec:rw-concepts}
\label{sec:rw-proto-recsys}     % prototype-landscape closing block (2026-07-14)
\vault{2026-06-11_1111_protomf-within-interpretable-ai}
% Positioning section (protocolled decision 2026-06-11_1111). Moved BEFORE
% ProtoMF at the merge (Matteo, 2026-07-14): teaches the vocabulary that lets
% ProtoMF then be read as an instance of interpretable-by-design.
% SUBSECTION STRUCTURE DELIBERATELY UNDECIDED (Matteo, 2026-07-14): the old
% three-subsection split (intrinsic-vs-post-hoc / x-to-c-to-y / additivity)
% is retired; Matteo decides the headers when working through the protocol
% vault entries. Content to be housed, as bullets:
\begin{itemize}
    \item [TBD] Intrinsic vs post-hoc explanation — Rudin stance; the nuance: post hoc concerns the \emph{naming}, not the model. \vault{2026-06-11_1220_posthoc-naming-vs-intrinsic-model-distinction, 2026-06-16_1808_intrinsic-over-posthoc-principle, 2026-06-08_1149_prototype-naming-and-similarity-routes, 2026-07-14_1522_industry-explanations-posthoc-templates-citable}
    \item [TBD] Prototypes and concepts, the x$\to$c$\to$y view: ProtoMF solves c$\to$y, leaves c-semantics open. \vault{2026-06-11_1642_protomf-solves-c-to-y-open-on-c-semantics}
    \item [TBD] Concept-side positioning: CBMs, SENN, attention-as-explanation; dc02's enforced grounding as contrast.
    \item [TBD] Additivity and attribution — the linearity$\leftrightarrow$interpretability argument (DIRECTIVE: becomes a full paragraph). \vault{2026-07-12_1200_linearity-interpretability-additive-attribution}
    \item [TBD] ``GAM-like in prototype space, not a GAM'' — one sentence. \vault{2026-06-11_1504_protomf-not-a-gam-additive-only-in-proto-space}
    \item [TBD] Marked slot: does the naming problem generalize across the prototype family? (standing research task)
    \item [TBD] Closing block: the prototype landscape as transition into ProtoMF — family survey $\to$ recsys neighborhood (survey level only; verdicts live in The Gap) $\to$ the member we build on. \vault{2026-06-16_2057_protomf-vs-ppm-positioning, 2026-07-11_1515_ppm-prominence-related-work-anchor, 2026-06-11_1522_part-based-prototypes-unlocked-by-features}
\end{itemize}

\section{ProtoMF: Prototype-Based Matrix Factorization}
\label{sec:bg-protomf}
\vault{2026-03-12_1640_phase2-6-paper-reread-complete}
\vault{2026-08-17_1937_base-settings-decision-shifted-bias-off-cosbias2x2-results, 2026-08-17_1128_shifted-cosine-baseline-asymmetry-popularity-channel}
% DIRECTIVE (Matteo, 2026-08-17): at the END of this section, bring up what
% is concealed in the shifted cosine (the structural popularity channel, the
% U/I side asymmetry) and how the cosbias2x2 shows it does not seem to impact
% much (metrics robust, I-side placebo, incumbent not dominated). The two
% annotated grid sheets carry the visual argument; (shifted, bias-off) is
% thereby evidenced, not inherited — final ratification with equal-treatment
% numbers at writing time (two-tier policy).
% The host family, taught at working depth. Mirrors the ProtoMF paper's own
% build order. Readable as an instance of interpretable-by-design (previous
% section supplies the vocabulary; its closing prototype-landscape block IS
% the transition into this section). PURE TEACHING since 2026-07-14 (second
% restructure): the old closing positioning subsection moved to the end of
% the IML section (survey half) and The Gap (verdict half) — this section
% only teaches the host, and ends on the regularizers' Part II tension,
% handing over to the mechanism axis (LightFM) next.

\subsection{User and Item Prototypes (U-ProtoMF, I-ProtoMF)}
\label{sec:bg-protomf-ui}
% shifted cosine similarity, prototype read-out, s^user/s^item.
% Decomposed from vault 2026-07-13_2205 (not prose as one blob).
\begin{itemize}
    \item Prototypes = anchors in the shared latent space; entities re-expressed by their relation to the anchors. \vault{2026-07-13_2205_anchors-combination-vs-membership-cluster-survives-in-naming}
    \item Combination, not membership: a graded affinity profile over all prototypes — multi-interest taste survives. \vault{2026-07-13_2205_anchors-combination-vs-membership-cluster-survives-in-naming}
    \item The profile IS the scoring representation: the score decomposes additively per prototype. \vault{2026-07-13_2205_anchors-combination-vs-membership-cluster-survives-in-naming}
    \item The cluster notion survives only in the \emph{naming} (centroid reading) — until the attributes arrive. \vault{2026-07-13_2205_anchors-combination-vs-membership-cluster-survives-in-naming}
    \item ACF: the anchor idea by literal name — lineage entry point.
\end{itemize}

\subsection{The Double-Tied Model (UI-ProtoMF)}
\label{sec:bg-protomf-tie}
\vault{2026-08-17_1823_correction-double-tie-projection-basis-and-ui-effective-count}
% weight sharing, compromised embeddings, tie trade-off (2026-03-18 entry).
% CORRECTION on record (2026-08-17_1823): at the double tie, only the ENTITY
% embeddings hold two jobs (cosine leg + linear leg); the projections are free
% W matrices and prototype norms stay pure gauge — teach the tie mechanics to
% this corrected statement, never the refuted "prototypes as projection basis".
\begin{itemize}
    \item A tied vector is a latent vector with two jobs — sharing constrains parameters, it does not change latent-vs-anchored. \vault{2026-03-18_1234_double-tie-explainability-tradeoff, 2026-07-13_2033_what-remains-latent-after-grounding-still-mf}
\end{itemize}

\subsection{Prototype Regularizers}
\label{sec:bg-protomf-reg}
% reg_batch / reg_proto geometric effect. NOT claimed as "pure"
% interpretability penalties anymore (Matteo, 2026-07-14): interpretability-
% motivated, but — like L1/L2 in NNs — they plausibly also act as generic
% regularization and affect accuracy; that is an empirical question.
\begin{itemize}
    \item The two regularizers shape reasoning \emph{legibility}, not transparency — transparency is architectural (the cosine read-out). \vault{2026-06-15_1238_protomf-two-regularizers-geometric-effect, 2026-06-11_1144_proto-regularizers-interpretability-not-accuracy}
    \item Interpretability-motivated, but plausibly dual-effect — like L1/L2, they may also move accuracy. \vault{2026-07-14_1152_outline-refinement-ch2-8-structural-decisions}
    \item [PLANNED] Small ablation: their effect on accuracy AND the legibility measures. \vault{2026-07-14_1152_outline-refinement-ch2-8-structural-decisions, 2026-08-17_1407_coverage-regularizer-uncollapses-hm-prototype-space}
    \item Hyperopt tunes them on accuracy alone — the tension Part~II picks up. \vault{2026-06-11_1144_proto-regularizers-interpretability-not-accuracy}
\end{itemize}

\section{Feature-Aware Matrix Factorization and LightFM}
\label{sec:bg-lightfm}
\label{sec:rw-feature-mf}   % old RW section folded in here (positioning pass)
\vault{2026-06-16_1715_feature-injection-design-fork, 2026-06-16_1824_feature-aware-baselines, 2026-06-16_1838_fm-gmf-suitability-hybrid-bridge}
% Axis 1 in one section: a SHORT positioning pass over the family, then
% LightFM taught at working depth — the mechanism our lineage builds on AND
% our CBF baseline.
% Positioning pass (one paragraph + citation cluster, old RW section): CMF,
% RLFM, SVDFeature, FM family (AFM/NFM/DeepFM/Wide&Deep), CTR/CDL, DropoutNet;
% "expressive but opaque" foil — features enter the SCORE, not the
% explanation vocabulary; grounding-bridge verdict (FM/GMF not THE mechanism);
% feature-injection fork: prediction-side vs prototype-side.
% Teaching pass (LightFM): entity = sum of feature embeddings (+ ID row);
% tags vs tags+ids arms; cold-start capability.
% EVALUATIVE PASSAGE (Matteo, 2026-07-14): did LightFM's composed embeddings
% buy any interpretability over basic MF — and if so, how? Verdict to argue:
% partial, and as a byproduct, not a goal.
%  - Gained: rows are ANCHORED, so the latent space becomes queryable —
%    Kula's own tag-similarity/analogy demonstrations; semantics carried by
%    features instead of IDs (which is also WHY cold-start works); in our
%    2.1 taxonomy: identity observed, coordinates fitted.
%  - NOT gained: no aggregated read-out surface. The bilinear score
%    decomposes only into feature-by-feature inner-product terms — quadratic,
%    unaggregated, not a usable explanation of a recommendation; nothing in
%    the model carries a printable name per component. The paper itself
%    motivates accuracy + cold-start, not interpretability.
%  - Payoff sentence: LightFM grounds the INGREDIENTS but has no
%    interpretable read-out to serve them to; ProtoMF has the read-out but
%    ungrounded ingredients — our lineage is the splice. Feeds The Gap
%    ("the gap is the combination, not an ingredient").
\begin{itemize}
    \item Collaborative = \emph{what trains} a model; content-based = \emph{how it is parameterized} — composed embeddings are content-parameterized, collaboratively trained. \vault{2026-07-13_2033_what-remains-latent-after-grounding-still-mf}
    \item The noid arms are that hybrid's extreme — why LightFM is the closest baseline, not decoration. \vault{2026-07-13_2033_what-remains-latent-after-grounding-still-mf}
\end{itemize}

\section{The Gap}
\label{sec:rw-gap}
\vault{2026-06-16_1808_intrinsic-over-posthoc-principle, 2026-07-13_1817_core-level-modality-light-positioning, 2026-07-14_1218_descriptor-gap-always-a-gap-fidelity-vs-legibility}
% Closing move (DBIS/VUB norm) — the pincer snaps shut. Structure (2026-07-14):
%  1. Close each axis with its failure mode: feature-aware MF is mature but
%     opaque (features in the score, not the explanation vocabulary);
%     prototype-based recommendation is interpretable-shaped, but every
%     prototype is named post hoc from its neighborhood — the naming stays
%     outside the model.
%  2. The checkable claim: no ProtoMF descendant grounds its prototypes in
%     catalogue attributes intrinsically (forward-citation sweep I1-I5; one
%     sentence each on why ProtoCF / UIPC-MF / ACF do not fill it).
%     SOLE OWNER of these verdicts since 2026-07-14: the closest-prior-art /
%     closest-competitor judgments live ONLY here — the IML section's
%     landscape block names those works at survey level without verdicts.
%  3. The gap is the COMBINATION, not an ingredient: composition exists
%     (LightFM), semantic anchoring exists in general ML (CBMs, annotated
%     prototypes) — nobody has put composition inside the prototype read-out
%     of a recommender.
%  4. Modality-light restatement (catalogue attributes only, vs review-corpus
%     methods; vault 2026-07-13_1817).
%  5. Hand-off: the host aimed at interpretability and stopped one step short
%     (echoes §1.2); this thesis closes that step with a staged lineage
%     (fI -> fU -> fUfI).
% OPTIONAL SHARPENER (marked slot): if the family-wide naming-problem
% research confirms, upgrade claim 2 to "post-hoc naming is the family-wide
% pattern; this is the first recsys member to break it."
